\subsection{The intrinsic weak class and single-square diagonals}\label{sec:B5T01}
A residual board retains the original coordinates of the ordinary queens board.
Its support $E$ has $k$ rows and $k$ columns. Each square belongs to
four lines, and two distinct lines intersect in at most one square.
Let $p$ be an interior Bethe stationary point with row and column masses equal to one.
For fixed constants $C,\gamma,\eta$, we define the deterministic class by
\[
 p_e\le C/k,\quad \operatorname{gap}_{RC}(p)\ge\gamma,
 \quad h_d=1-m_d\ge\eta,\quad
 m_d=\sum_{e\in d}p_e\le\varepsilon\le\varepsilon_0.
\]
The RC gauge is $(\mathbf1_R,-\mathbf1_C)$, and the gap is taken on its orthogonal complement. We also require
RC connectivity. For $k\ge k_0$ we have $p<1/4$. Unlike the macroscopic class, this class does not
require $k\ge\delta_0n$, and the constants do not depend on the number of diagonals. The bound on $p$
forces every row and every column to have at least $k/C$ squares. The overall density of the mask
alone cannot replace these conditions.

The original geometry remains part of the problem. The intrinsic class still takes account
of which squares share a diagonal.

\paragraph{Discrete lower bound on the mass of a diagonal.}
Since each row sums to one and has at most $k$ squares, we have
\[
 \max_e p_e\ge1/k,\qquad \max_d m_d\ge1/k.
\]
In particular, a bound of the form $\ell_{\max}\le Ak/n$ would require $k^2\ge n/A$.

\paragraph{Exact elimination of a single-square diagonal.}
Write $q=1-p$ and allow a multiplicity $d\in\{2,3,4\}$, so that
\[
 \tau_d=\frac{pq}{1-dp},\qquad b_d=\frac{1-dp}{q^d}.
\]
The single-square line has vacancy $q$. Its pivot and Schur complement are
\[
 q+\tau_d=q\frac{1-(d-1)p}{1-dp},\quad
 \tau_d-\frac{\tau_d^2}{q+\tau_d}=\tau_{d-1},\quad
 b_d(q+\tau_d)=b_{d-1}.
\]
The product of the local factor and the determinant is therefore preserved exactly.
The Bernoulli entropy of this factor cancels the change in multiplicity
in the Bethe entropy. The RC gauge is unchanged and has zero coordinate
on the eliminated line. Two single-square lines take us from $d=4$ to $d=2$.
A line with no squares contributes a pivot equal to one.

The matrix row cannot be deleted while retaining $\tau_4,b_4$. To avoid having to track
variable multiplicities anew, Sections~\ref{sec:B5T02}--\ref{sec:B5T05} retain all four
incidences of each square, even for single-square lines. The number of such lines does
not affect the mass bounds. The rook model can be represented by adding two
single-square diagonals to each square and eliminating them with this identity.

\paragraph{Scope.}
The predictor and the RC cofactor are those of Sections~\ref{sec:B3T01} and~\ref{sec:B4T01}.
We make no claim about typical membership in this class or the completability of its members.
